\chapter{Cystic Fibrosis}

\section*{Definition and Overview}
Cystic fibrosis (CF) is an inherited condition in which a faulty gene causes the body to produce thick, sticky mucus. This mucus builds up mainly in the lungs and the digestive system, where it clogs the airways, traps bacteria, and blocks the release of digestive enzymes from the pancreas. CF is present from birth, but it may be diagnosed at different ages depending on how severe the symptoms are and whether newborn screening is used. Although there is no cure, treatments that keep mucus thin, clear airways, and maintain good nutrition have greatly improved life expectancy.

\section*{Epidemiology}
Cystic fibrosis is one of the more common inherited conditions in people of European ancestry, affecting roughly 1 in 2,500 to 3,500 newborns in these populations. It is less common in people of African, Asian, and Middle Eastern descent. Both sexes are affected equally.

The condition requires a faulty copy of the gene from each parent, and around 1 in 25 people of European descent is a carrier without symptoms. Diagnosis most often occurs in childhood, but a small number of people are diagnosed in adulthood, usually with milder disease. In countries with newborn screening, most cases are now picked up in the first weeks of life.

\section*{Aetiology and Pathophysiology}
\begin{itemize}
    \item \textbf{Genetic cause:} caused by mutations in the \textit{CFTR} gene, which normally produces a protein that acts as a channel for salt and water in and out of cells
    \item \textbf{Inheritance:} autosomal recessive -- a child must inherit a faulty copy from both parents to develop CF
    \item \textbf{Abnormal mucus:} in CF the channel does not work properly, so salt and water movement is disrupted and mucus becomes abnormally thick and sticky
    \item \textbf{Lung damage:} sticky mucus blocks small airways and provides a breeding ground for bacteria, leading to repeated infection and progressive inflammation
    \item \textbf{Pancreatic problems:} thickened secretions block pancreatic ducts, so digestive enzymes cannot reach the intestine and food is poorly absorbed
    \item \textbf{Other organs:} similar secretions affect sweat glands (very salty sweat), the liver, bile ducts, and the reproductive system
\end{itemize}

\section*{Clinical Features}
\begin{itemize}
    \item Persistent, productive cough with thick sputum
    \item Frequent chest infections, wheeze, and breathlessness
    \item Poor weight gain and growth despite a good appetite (malabsorption)
    \item Greasy, bulky, foul-smelling stools
    \item Very salty-tasting skin
    \item Salty sweat test results (elevated chloride)
    \item Clubbing of the fingers in established disease
    \item Blocked nasal passages, sinus problems, and nasal polyps
    \item In males, absence of the vas deferens causing infertility
    \item Delayed puberty in some children with poor nutrition
\end{itemize}

\section*{Differential Diagnosis}
\begin{itemize}
    \item Asthma or recurrent wheeze
    \item Primary ciliary dyskinesia (a condition causing impaired clearance of mucus)
    \item Recurrent pneumonia or bronchiectasis from other causes
    \item Coeliac disease or other causes of malabsorption and poor growth
    \item Allergic bronchopulmonary aspergillosis
    \item Immunodeficiency syndromes with recurrent infections
    \item Chronic lung disease of prematurity or chronic aspiration
\end{itemize}

\section*{When to Seek Medical Care}
See a doctor promptly if a child or adult with CF has increasing cough, new breathlessness, a change in the colour or amount of sputum, unexplained fever, or poor appetite. Anyone without a diagnosis who has a persistent cough, frequent chest infections, salty skin, or poor weight gain should also be assessed.

\textbf{Call 000 (or local emergency services) for:}
\begin{itemize}
    \item Severe difficulty breathing or rapid breathing
    \item Coughing up blood (haemoptysis)
    \item Sudden severe chest pain, especially with breathlessness (possible collapsed lung)
    \item Confusion, extreme drowsiness, or blue-tinged lips or skin
\end{itemize}

\section*{Diagnosis and Investigations}
\begin{itemize}
    \item \textbf{Sweat test:} the standard diagnostic test, measuring chloride levels in sweat; high levels support the diagnosis
    \item \textbf{Genetic testing:} identifies \textit{CFTR} gene mutations and is used to confirm the diagnosis and guide treatment
    \item \textbf{Newborn screening:} a blood test in the first days of life detects elevated immunoreactive trypsinogen, prompting further testing
    \item \textbf{Nasal potential difference or intestinal current measurement:} specialised tests used when the diagnosis is unclear
    \item \textbf{Chest imaging:} chest X-ray or CT to assess lung changes such as bronchiectasis
    \item \textbf{Lung function tests:} spirometry to monitor airflow obstruction over time
    \item \textbf{Sputum culture:} identifies the bacteria present to guide antibiotics
    \item \textbf{Blood tests:} pancreatic enzyme levels, vitamin levels, and liver function
\end{itemize}

\section*{Management}
\begin{itemize}
    \item \textbf{Chest physiotherapy:} daily airway clearance techniques to loosen and remove mucus from the lungs
    \item \textbf{Inhaled medicines:} bronchodilators, hypertonic saline, and DNA-dissolving agents (such as dornase alfa) to thin mucus and open airways
    \item \textbf{Antibiotics:} oral, inhaled, or intravenous courses to treat and prevent lung infections
    \item \textbf{CFTR modulator therapies:} newer medicines that correct the faulty protein's function in people with eligible gene types
    \item \textbf{Pancreatic enzyme supplements:} taken with food to allow proper digestion and absorption
    \item \textbf{Nutritional support:} a high-energy, high-fat diet with fat-soluble vitamin supplements and regular dietitian review
    \item \textbf{Exercise:} regular physical activity improves airway clearance and fitness
    \item \textbf{Surgery:} occasionally needed for nasal polyps, sinus disease, or advanced lung disease (lung transplantation)
    \item \textbf{Specialist care:} care in a multidisciplinary CF centre with regular review
\end{itemize}

\section*{Complications}
\begin{itemize}
    \item Chronic lung infections, especially with \textit{Pseudomonas aeruginosa} and other resistant bacteria
    \item Bronchiectasis and progressive decline in lung function
    \item Coughing up blood or major haemorrhage from damaged airways
    \item Pneumothorax (collapsed lung)
    \item Respiratory failure in advanced disease
    \item Diabetes related to CF, from scarring of the pancreas
    \item Liver disease and cirrhosis from thickened bile
    \item Biliary obstruction and gallstones
    \item Malnutrition, poor growth, and delayed puberty
    \item Infertility in males, and reduced fertility in some females
    \item Osteoporosis and low bone density
    \item Psychological and social impact of a lifelong condition
\end{itemize}

\section*{Prognosis and Outcomes}
Life expectancy for people with cystic fibrosis has increased dramatically over recent decades. In countries with modern care, median survival now exceeds the fifth decade of life, and many people live well beyond that. Outcomes depend on the specific gene mutations, how early treatment starts, adherence to daily therapy, and access to specialist care. CFTR modulator therapies have substantially improved lung function and quality of life for eligible people. The condition remains progressive, and some people eventually develop respiratory failure requiring consideration of lung transplantation. Individual prognosis should be discussed with the treating CF team.

\section*{Prevention and Self-Care}
\begin{itemize}
    \item Newborn screening allows early diagnosis and treatment before major damage occurs
    \item Genetic counselling and carrier testing for people with a family history of CF, or whose partner is a carrier
    \item Daily airway clearance and inhaled treatments as prescribed
    \item Take pancreatic enzymes with every meal and snack as directed
    \item Keep up to date with vaccinations, including influenza and pneumococcal vaccines
    \item Maintain good hand hygiene and avoid contact with people who are unwell
    \item Follow a high-energy, nutritious diet and attend regular dietitian reviews
    \item Stay physically active and avoid tobacco smoke and air pollution
    \item Attend regular specialist clinic reviews to monitor lung function, nutrition, and growth
\end{itemize}

\section*{Sources and Further Reading}
The following organisations provide reliable information on cystic fibrosis and its management.

\begin{itemize}
    \item National Health Service. \textit{NHS conditions guide on cystic fibrosis.} https://www.nhs.uk/conditions/
    \item Mayo Clinic. \textit{Cystic fibrosis: symptoms, causes, and treatment.} https://www.mayoclinic.org/
    \item Centers for Disease Control and Prevention. \textit{Cystic fibrosis and newborn screening information.} https://www.cdc.gov/
    \item World Health Organization. \textit{Global guidance on chronic respiratory disease.} https://www.who.int/
    \item MedlinePlus. \textit{Consumer health information on cystic fibrosis.} https://medlineplus.gov/
    \item MSD Manual. \textit{Cystic fibrosis -- professional overview.} https://www.msdmanuals.com/
\end{itemize}
